%% kapa: exam=2 round=20 year=2009 period=1 subject=appraisal_practice subject_ko=감정평가실무 seq=2 number=2 points=25
%% source: Q-Net 감정평가사 기출문제 (cst00309) artlSeq=5024180 file=extracted/1008/5024180/2009년 제20회 감정평가사 2차시험 기출문제/1교시.hwp
%% conversion: kordoc 4.24.51 -> tools/gen2.py (auto); review=proofread-sonnet; flags=none
\begin{longquestion}{【문제 2】}
투자가 K씨는 다음의 부동산(A, B)중 하나에 투자하려고 한다. 감정평가사인 Y씨는 대상토지가치의 타당성 검토를 의뢰받았다. 동일한 금액을 투자할 경우 적절한 투자방안을 결정하고 그 이유를 설명하시오. (25점) \par
<자료 1> \par
1) 대상부동산 A \par
① 소재지: ○○시 ○○구 ○○동 ○○번지 \par
② 지목: 대 \par
③ 면적: 500m\textsuperscript{2} \par
④ 이용상황: 나지 \par
⑤ 도시계획: 일반상업지역 \par
2) 대상부동산 B \par
① 소재지: △△시 △△구 △△동 △△번지 \par
② 지목: 대 \par
③ 면적: 1,000m\textsuperscript{2} \par
④ 이용상황: 나지 \par
⑤ 도시계획: 일반상업지역 \par
<자료 2> A 투자계획 \par
1) 대상부동산 A에 업무용 건물을 건축하여 임대할 예정으로 연면적 2,000m\textsuperscript{2}, 지상 5층(각층 동일면적, 지하층 없음)으로 계획 중임. \par
2) 건물신축비용은 950,000원/m\textsuperscript{2}임. \par
3) 예상 지불임대료 \par
① 1층: 월 임대료 30,000원/m\textsuperscript{2}, 월 관리비 9,000원/m\textsuperscript{2}로 책정됨. \par
② 2, 3층: 월 임대료 15,000원/m\textsuperscript{2}, 월 관리비 9,000원/m\textsuperscript{2}로 책정됨. \par
③ 4, 5층: 월 임대료 12,000원/m\textsuperscript{2}, 월 관리비 9,000원/m\textsuperscript{2}로 책정됨. \par
4) 보증금은 월 임대료의 12개월분으로 함. (관리비는 포함하지 않음) \par
5) 유사부동산들의 수익률 범위를 조사한 바 그 중 70\%는 수익률 12\%, 15\%는 수익률 13\%, 15\%는 수익률 11\%로 조사됨. \par
<자료 3> B 투자계획 \par
1) 대상부동산 B에 할인점을 건축하여 수수료매장으로 운영할 예정임. 연면적 2,000m\textsuperscript{2}, 지상 2층(지하층 없음)으로 계획 중임. \par
2) 건물신축비용은 700,000원/m\textsuperscript{2}임. \par
3) 할인점은 수수료매장으로 운영하여 매출액의 2\%를 지불임대료로 받음. △△시에 거주하는 가구 수의 40\%가 대상할인점을 이용할 것으로 예상되며 △△시는 가구당 평균인구수가 3.5인, 가구당 연간 평균소득은 30,000,000원임. 대상할인점을 이용하는 가구는 평균적으로 소득의 3\%를 대상할인점을 통해 물품을 매입할 것으로 조사되었음. \par
4) 별도의 관리비는 징수하지 아니함. \par
5) 보증금은 월 지불임대료의 12개월분으로 함. \par
6) 유사부동산들의 수익률범위를 조사한 바 그 중 70\%는 수익률 12\%, 15\%는 수익률 14\%, 15\%는 수익률 10\%로 조사됨. \par
<자료 4> 기타 \par
1) ○○시 인구규모: 1,000,000명이며 미미하게 증가추세를 보임. \par
△△시 인구규모: 300,000명으로 정체 중임. \par
2) 보증금운용이율은 8\%로 적용하며 시장의 무위험이자율은 7\%로 적용함. \par
3) 필요제경비는 대상부동산 A의 경우 관리비 수령액의 80\%로 하며, 대상부동산 B의 경우 지불임대료의 30\%로 함. \par
4) 임대가능성, 대손, 공실 등은 고려하지 아니함. \par
5) 환원이율은 무위험이자율에 위험률을 합한 율로 적용하며 위험률의 산정은 유사부동산 수익률범위의 표준편차를 적용함. \par
\end{longquestion}
